\section{Reducing the anchor coefficients and shortening its constant chain}
\label{sec:p3-affine-anchor-refinement}

The following refinements are exact transformations of the supplied
quadratic residual systems.  They require no additional arithmetic
theorem.  The original tables remain available; the transformed
tables are separate artifacts with certificates referring to their
frozen parents.

\subsection{An affine change in the three private anchor coordinates}

Put
\[
 \kappa=3^{2^{28}},\qquad s=3^{2^{27}},\qquad s^2=\kappa.
\]
In the $1315$-coordinate table the anchor residual is
\[
 R=-3A^2+16u^2+96v^2-32w^2+32\kappa,
 \qquad A=X_{52}+X_{53}-X_{54},
\]
where $(u,v,w)=(X_{55},X_{56},X_{57})$.  These three coordinates
occur in no other residual.  The other three displayed coordinates
are constrained arithmetic values, so they cannot be treated as
independent anchor witnesses.

Make the rational affine substitution
\begin{equation}\label{eq:p3-anchor-affine-map}
 u=s+\frac{U-3V+3W}{4},\qquad
 v=\frac{s}{2}+\frac V4,\qquad
 w=\frac{3s}{2}+\frac W4.
\end{equation}
Its linear determinant is $1/64$, and its inverse is
\[
 U=4u+12v-12w+8s,\qquad V=4v-2s,\qquad W=4w-6s.
\]
It is therefore a bijection over every characteristic-zero field,
fixing the input and every other coordinate.  Direct expansion gives
\begin{equation}\label{eq:p3-translated-anchor}
 \widetilde R=-3A^2+U^2-6UV+6UW+15V^2-18VW+7W^2+8sU.
\end{equation}
Indeed the constant terms sum to
$16s^2+24s^2-72s^2+32s^2=0$, and the linear terms sum to $8sU$.
All coefficients of the new residual are integers, with $s$ regarded
as the specified fixed integer, not a coordinate.

\begin{proposition}\label{prop:p3-translated-quartic}
Replacing $R$ by $\widetilde R$ in the frozen $1315$-coordinate
residual list gives a polynomial $\widetilde P_{1315}$ with exactly
the same input projection as $P_{1315}$ over every characteristic-zero
field.  Its degree is exactly four, its primitive integer coefficient
height is
\[
                  H(\widetilde P_{1315})=64\kappa
                                      <2^{429496736},
\]
and its complete table has $29021$ nonzero ordinary coordinate
monomials.  It has $1272$ quadratic residuals and $1314$ existential
witnesses.  In particular, under the already stated arithmetic inputs,
it defines $\mathbb Z$ in $\mathbb Q$.
\end{proposition}
\begin{proof}
The polynomial identity
\[
 \widetilde P_{1315}
       =P_{1315}\circ\text{the substitution in
                           \eqref{eq:p3-anchor-affine-map}}
\]
and the explicit inverse prove equality of input projections.
This conclusion concerns the literal single polynomial as well as
the common residual zero set; it does not use positivity of squares.
No new coordinate or equation is introduced, and every residual
remains quadratic.

For the height assertion write
\[
 B=-3A^2+U^2-6UV+6UW+15V^2-18VW+7W^2.
\]
Then $\widetilde R^2=B^2+16sUB+64s^2U^2$.  Since $U,V,W$
occur nowhere else, the coefficient of the monomial $U^2$ in
the sum of squares is exactly $64s^2$.  It is the only coefficient
involving $s^2$.  The coefficients involving $s$ have absolute value
at most $288s$.  The independently expanded table records that
every remaining integer coefficient has absolute value less than
$2^{1020}$, and no monomial receives two different powers of $s$.
The value $s=3^{2^{27}}$ makes $64s^2$ larger than both other bounds.
The input coefficient of $X^4$ is exactly one.  Thus the integer
content is one, the height is exactly $64s^2=64\kappa$, and the
degree is exactly four.  Finally $3^5<2^8$ gives
\[
 \log_2(64\kappa)<6+\frac85\,2^{28}<429496736.
\]
The displayed residual list and its exact convolution give the
stated monomial count.
\end{proof}

\subsection{Eliminating the last coordinate in the constant chain}

In the ordinary-coefficient $1339$-coordinate table, the coordinate
$X_{81}$ occurs in exactly two residuals:
\[
                   X_{80}^2-X_{81},\qquad R_0+32X_{81},
\]
where
\[
 R_0=-3(X_{52}+X_{53}-X_{54})^2
                       +16X_{55}^2+96X_{56}^2-32X_{57}^2.
\]
Substitute $X_{81}=X_{80}^2$, delete its defining residual, and
renumber all subsequent coordinates by subtracting one.  The second
residual becomes $R_0+32X_{80}^2$ and is still quadratic.

\begin{proposition}\label{prop:p3-small-coefficient-1338}
The resulting system has $1338$ total coordinates and $1295$
quadratic residuals.  Its sum of squares $P_{1338}$ has degree exactly
four, $29059$ nonzero ordinary integer monomials, primitive integer
content one, and coefficient height less than $2^{1020}$.
Its rational input projection equals that of $P_{1339}$.
\end{proposition}
\begin{proof}
On a common zero of the parent residuals, the first displayed equation
uniquely determines $X_{81}=X_{80}^2$.  Deleting this coordinate
therefore gives a solution of the new system.  Conversely, inserting
that square in a new solution restores all parent equations.  This
is a bijection of common residual zero sets over every
characteristic-zero field and fixes the input.

Over a formally real field, in particular over $\mathbb Q$, a sum
of squares vanishes precisely when all its residuals vanish.  Hence
the same equivalence holds for the literal sum-of-squares input
projections over those fields.  No such converse for the literal
polynomials over nonreal fields is inferred from this elimination.

The substitution keeps every residual quadratic.  Exact convolution
of the remaining residuals gives $29059$ nonzero coefficients, all
of absolute value below $2^{1020}$.  The unchanged coefficient of
the input monomial $X^4$ is one, proving content one and degree four.
The complete maximum coefficient is recorded in the construction
report.  The coordinate and residual counts each decrease by one.
\end{proof}

\subsection{Inlining the three initial constants}

The first three coordinates of the remaining constant chain are
uniquely determined as
\[
                  X_{58}=3^{32},\qquad X_{59}=3^{64},
                  \qquad X_{60}=3^{128}.
\]
Their three defining residuals are
\[
       3^{32}-X_{58},\qquad X_{58}^2-X_{59},
       \qquad X_{59}^2-X_{60}.
\]
They have no other occurrences except in the next equation of
the chain.  Substitute the three constants and delete these
coordinates and residuals.  The first remaining chain equation
becomes $3^{256}-X_{61}=0$ before renumbering, or
$3^{256}-X_{58}=0$ afterwards.

\begin{theorem}[The preferred ordinary-coefficient quartic]
\label{thm:p3-small-coefficient-1335}
The supplied polynomial $P_{1335}$ has $1335$ total coordinates,
$1334$ existential witnesses, $1292$ quadratic residuals, degree
exactly four, and $29050$ nonzero ordinary integer monomials.
Its primitive coefficient height is unchanged from $P_{1338}$ and
is less than $2^{1020}$.  Under the same arithmetic hypotheses,
it defines $\mathbb Z$ in $\mathbb Q$.
\end{theorem}
\begin{proof}
Each eliminated coordinate equals the displayed fixed integer
at every common zero of the parent residuals.  Substitution and
reinsertion are inverse maps on the corresponding common zero
sets, over every characteristic-zero field, and fix the input.
All remaining residuals have degree at most two.  Over $\mathbb Q$
their sum of squares therefore has the same input projection.

The first remaining chain residual contributes the constant
$3^{512}$ on squaring, and
\[
                             3^{512}<2^{812}<2^{1020}.
\]
For an elementary check of the first inequality, $3^{41}<2^{65}$
and $512\cdot65<812\cdot41$.  The only changed constant
coefficient is the old one minus $3^{64}$ plus $3^{512}$; it
has bit length at most $826$.  The new linear coefficient has
absolute value $2\cdot3^{256}<2^{407}$.  Every other coefficient
is inherited, up to coordinate renumbering, or removed.  The
previous coefficient of maximum height occurs in an unaffected
monomial, so the exact height is unchanged.  These assertions
are also checked against the complete coefficient table.

Twenty constant-chain coordinates remain, representing
$3^{2^8},3^{2^9},\ldots,3^{2^{27}}$.  The final square is
used directly in the anchor.  Thus the ordinary-coefficient
presentation adds twenty coordinates to the $1315$-coordinate
literal-constant presentation.  Exact convolution gives $29050$
nonzero terms.  The unchanged coefficient of the input monomial
$X^4$ is one, proving primitive content one and degree four.
\end{proof}

The preferred fresh packages are
\nolinkurl{research/phase3/quartic/quartic-direct-translated/} and
\nolinkurl{research/phase3/quartic/quartic-direct-small-coefficients-1335/}.
The intermediate $1338$-coordinate package is retained separately
under \nolinkurl{quartic-direct-small-coefficients-1338/}.
Each contains the full residual list, complete coefficient table,
parent hash, and exact transformation certificate.  Their producer
is \nolinkurl{transform_direct_quartics.py}.  It verifies both the
residual substitutions and the full polynomial identities against
the frozen parent tables.  The separate program
\nolinkurl{verify_final_transforms_independent.py} imports no producer
code and independently checks every transformed residual, every
expanded coefficient, all coordinate maps, and the stated heights
and supports.  All three checks pass.  The translated table uses
the coefficient symbol $S$ for the fixed integer $s$; an entry
$[\mathbf m,k,c]$ means $cs^k\prod_{i\in\mathbf m}X_i$.
This symbol contributes no variable and no degree in the quantified
coordinates.
